\documentclass[10pt,a4paper]{amsart}
\usepackage[margin=19mm]{geometry}
\usepackage{amsmath,amssymb,mathtools}
\usepackage[scr=rsfs,cal=euler]{mathalfa}
\usepackage{xfrac}
\usepackage[unicode,hidelinks,bookmarks=false]{hyperref}
\newcommand{\ie}{i.e.\ }
\newcommand{\cf}{cf.\ }
\newcommand{\resp}{respectively\ }
\newcommand{\Subsection}{\subsection}
\providecommand{\nobreakdash}{\nobreak\hbox{-}}
\setlength{\emergencystretch}{2em}
\multlinegap0pt
\usepackage{mathtools}
\usepackage{amssymb}
\usepackage{xcolor}
\usepackage{tikz}
\usepackage{enumerate}
\newtheorem{proposition}{Proposition}
\newtheorem{theorem}{Theorem}
\newcommand{\ZZ}{\mathbb Z}
\newcommand{\ba}{\mathbf{a}}
\newcommand{\bc}{\mathbf{c}}
\newcommand{\be}{\mathbf{e}}
\newcommand{\cF}{\mathcal{F}}    % spanning forests set
\newcommand{\eps}{\varepsilon}
\newcommand{\fC}{\mathfrak C}    % parking complex +
\newcommand{\kp}{\varkappa}
\renewcommand{\mathcal}{\mathscr}
\title{The Loopy Polynomial: from Tutte's Universal $V$-Function to Bizonotopal Geometry}
\author{Anatol Kirillov, Gleb Nenashev, Boris Shapiro, Arkady Vaintrob}
\date{Source: arXiv:2609.07728v1 (CC BY 4.0)}
\begin{document}
\maketitle

\noindent\emph{Source and licence.} The Loopy Polynomial: from Tutte's Universal $V$-Function to Bizonotopal Geometry. Version arXiv:2609.07728v1 (CC BY 4.0), distributed by arXiv under the Creative Commons Attribution 4.0 International licence, http://creativecommons.org/licenses/by/4.0/, as recorded in the arXiv licence metadata for this exact version and shown on the abstract page https://arxiv.org/abs/2609.07728v1. Full licence notice: \texttt{licenses/arxiv-2609.07728-CC-BY-4.0.txt}.\par\smallskip

\noindent\textit{Editorial scope.} Counted unit: the article's theorem thm:polyhedral-parking-cells (source0.tex lines 5180--5240). The article's complete original proof of it is delivered with it (source0.tex lines 5242--5313, one \texttt{proof} environment of 72 lines). No mathematics is written, repaired or re-derived: the statement and the proof are the article's own lines, in the article's own notation, and no retained line is substituted or re-flowed.  Retained uncounted as the source-local context the proof needs: lines 4933--4973; lines 5000--5007; lines 5176--5178.  Nothing was omitted from any retained span, so the package carries no omission record.  No retained line contains a citation, so the wrapper carries no bibliography and no bibliography entry.  No retained line contains a figure environment, an image-inclusion command or drawing code, so no image is delivered and the package declares no asset dependency.  Editorial formatting only, in addition: an AMS \texttt{amsart} preamble that restates the article's own declarations and macros used by the retained lines, namely the environments \texttt{proposition}, \texttt{theorem} on separate counters, together with the standard packages those lines need.  The retained environments print in this document as Proposition 1, Theorem 1 in order of appearance, whereas the article prints the same items with the numbering of its own full document.  The complete native source of the article as distributed in the arXiv e-print for this exact version is bundled unchanged as \texttt{sources/209655/source0.tex}.
\begin{proposition}[Deletion-contraction geometry]
\label{prop:polyhedral-section}
Let $e=uv$ be a non-loop edge oriented from $u$ to $v$.
\begin{enumerate}
\item[(i)] $\Gamma_e(G)$ is the support of a polyhedral subcomplex of
$\partial P_G$.  More explicitly,
\begin{equation}
 \Gamma_e(G)=
 \{\ba\in P_G:a_u=0\}
 \;\cup\!
 \bigcup_{\substack{B\subseteq V\\v\in B,\ u\notin B}}
 \{\ba\in P_G:\ba(B)=\kp_G(B)\}.
\label{eq:upper-boundary-faces}
\end{equation}

\item[(ii)] The restriction
\[
   \pi|_{\Gamma_e(G)}:\Gamma_e(G)\longrightarrow P_{G/e}
\]
is a piecewise integral-affine homeomorphism. Its inverse is the maximal-lift
section $s_e(\bc)=\ba(U(\bc))$.  In particular,
\begin{equation}
   \Gamma_e(G)\cap\ZZ^V
   =s_e(P_{G/e}\cap\ZZ^{V'})=A_e.
\label{eq:gamma-lattice-points}
\end{equation}

\item[(iii)] The deletion polytope has the intrinsic description
\begin{equation}
   D_e(G)=\{\ba\in P_G:\ba+d_e\in P_G\}.
\label{eq:deletion-movable}
\end{equation}
Therefore $D_e(G)\cap\Gamma_e(G)=\varnothing$ and  $D_e(G)\cap\ZZ^V=B_e.$

\item[(iv)] The residual region
  \[
   R_e(G):=P_G\setminus\bigl(D_e(G)\cup\Gamma_e(G)\bigr)
  \]
contains no lattice points.
\end{enumerate}
\end{proposition}

\begin{equation}
 U(\bc)=\min\!\left(
 c_w,
 \min_{\substack{B\subseteq V\\v\in B,\ u\notin B}}
 \bigl(\kp_G(B)-\bc(B\setminus\{v\})\bigr)
 \right)
\label{eq:U-piecewise-linear}
\end{equation}

\begin{equation}  \label{eq:cells}
  C_F:=\fC_F\cap\ZZ^V.
\end{equation}

\begin{theorem}
  [Polyhedral parking cells]
  \label{thm:polyhedral-parking-cells} Fix an ordering and an orientation of the edges of $G$.
For every spanning
forest $F$, the geometric parking complex $\fC_F\subseteq P_G$ has the following properties.
\begin{enumerate}
\item[(i)] The map
\[
   \Phi_F:\Box_F\longrightarrow\fC_F
\]
is a piecewise-linear homeomorphism.  After subdividing $\Box_F$ by
the finitely many domains of linearity of the maximal-lift sections occurring  along the path,
$\fC_F$ is the support of a finite rational polyhedral complex.
On every cell, $\Phi_F$ and its inverse are integral-affine.

\item[(ii)] If $F\ne F'$, then
\begin{equation}
   \fC_F\cap\fC_{F'}=\varnothing.
\label{eq:geometric-cells-disjoint}
\end{equation}

\item[(iii)] The lattice points of the geometric parking complexes give a partition into parking cells~\eqref{eq:cells}
\begin{equation}
   P_G\cap\ZZ^V
   =\bigsqcup_{F\in\cF(G)}C_F.
\label{eq:parking-lattice-partition}
\end{equation}
Moreover $\Phi_F$ restricts to a lattice bijection
\begin{equation}
   \prod_{T\in c(F)}
      \{0,1,\ldots,|E(T)|+\eps_F(T)\}
   \xrightarrow{\ \sim\ } C_F.
\label{eq:parking-chain-product}
\end{equation}

\item[(iv)] Put
\begin{equation}
   d(F):=|E(G)|-|F|-\eps(F),
   \qquad
   \eps(F)=\sum_{T\in c(F)}\eps_F(T).
\label{eq:parking-shift}
\end{equation}
Then for every $x=(x_T)_{T\in c(F)}\in\Box_F$ we have
\begin{equation}
   \sum_{v\in V}(\Phi_F(x))_v  = d(F)+\sum_{T\in c(F)}x_T.
\label{eq:parking-grading}
\end{equation}
Therefore
\begin{equation}
 \sum_{\ba\in C_F}q^{|\ba|} = q^{d(F)} \prod_{T\in c(F)}\left(1+q+\cdots+ q^{|E(T)|+\eps_F(T)}\right).
\label{eq:parking-cell-enumerator}
\end{equation}

\item[(v)] The complement
\begin{equation}
   P_G\setminus\bigcup_{F\in\cF(G)}\fC_F
\label{eq:lattice-free-complement}
\end{equation}
is lattice-free.
\end{enumerate}
\end{theorem}

\begin{proof}
The terminal polytope at the leaf corresponding to $F$ is exactly the box
$\Box_F$.  Reversing the leaf path, every deletion step applies an injective
integral-affine translation, while every contraction applies the injective maximal-lift section of Proposition~\ref{prop:polyhedral-section}.  Thus their
composition $\Phi_F$ is continuous and injective.

Each maximal-lift section is piecewise integral-affine because its upper coordinate
is the minimum of finitely many integral-affine functions, as in~\eqref{eq:U-piecewise-linear}.
Subdivide $\Box_F$ by pulling back all the domains of linearity that occur along the path
and taking their common refinement.
This gives a finite rational polyhedral subdivision on which $\Phi_F$ is integral-affine cell by cell.
Because $\Phi_F$ is globally injective, for any two cells $C_1,C_2$ of this subdivision one has
\[
   \Phi_F(C_1)\cap\Phi_F(C_2)=\Phi_F(C_1\cap C_2).
\]
The intersection $C_1\cap C_2$ is a common face of $C_1$ and $C_2$, and the restrictions
of $\Phi_F$ to the cells are affine bijections.
Hence its image is a common face of $\Phi_F(C_1)$ and $\Phi_F(C_2)$.
The image cells therefore form a finite rational polyhedral complex.
The inverse at each step is either a translation by $-\be_u$ or the integral linear projection $\pi$.
Therefore the inverse on every image cell is integral-affine as well.  This proves (i).

For (ii), take distinct forests $F$ and $F'$ and compare their leaf paths.
Let the first difference occur when a non-loop edge $e$ is processed in some intermediate graph $H$.
One path takes deletion, so at that stage its image is contained in
\[
   D_e(H)=\be_u+P_{H-e},
\]
whereas the other takes contraction, so its image is contained in
$\Gamma_e(H)$.  These two sets are disjoint by
Proposition~\ref{prop:polyhedral-section}(iii).  Up to that stage the two leaf
paths coincide, so the sequence of maps already applied is the same for both.
Therefore the remaining maps carrying the two images back to $P_G$ are
literally the same composition, and it is injective.
An injective map carries disjoint sets to disjoint sets, so the final images $\fC_F$ and $\fC_{F'}$ remain disjoint.

For (iii), every map used in the construction respects the relevant integer
lattices in both directions.  This is immediate for translations.  For a
maximal-lift section it follows from
Proposition~\ref{prop:polyhedral-section}(ii): an integral point of the
contracted polytope has an integral maximal lift, and the inverse projection
sends integral points to integral points.  Hence
\[
   \fC_F\cap\ZZ^V
   =\Phi_F(\Box_F\cap\ZZ^{c(F)}),
\]
and the lattice points of $\Box_F$ are exactly the product in
\eqref{eq:parking-chain-product}.

It remains to see that every lattice point of $P_G$ occurs.  At one
non-loop edge, Proposition~\ref{prop:polyhedral-section}(iii)--(iv) says that
every lattice point belongs to exactly one of the deletion image $D_e(H)$ and
the contraction boundary $\Gamma_e(H)$.  Repeating this decision down the
deletion-contraction tree eventually sends every lattice point to a unique
loop-only leaf.  The leaf is indexed by a unique spanning forest, proving
\eqref{eq:parking-lattice-partition}.

For (iv), a contraction section preserves total coordinate sum, whereas every
genuine deletion translation raises it by one.  Along the leaf associated with
$F$, the contracted edges are the $|F|$ forest edges, and the nonforest edges
that have become loops are precisely the $\eps(F)$ externally active
edges.  Hence the number of genuine deletions is
\[
   |E(G)|-|F|-\eps(F)=d(F).
\]
At the terminal box the total coordinate sum is $\sum_Tx_T$, which proves
\eqref{eq:parking-grading}.  Restricting to the lattice points of the box gives
\eqref{eq:parking-cell-enumerator}.

Finally, (v) follows immediately from (iii): every lattice point of $P_G$ lies
in one of the $\fC_F$, so their complement contains none.
\end{proof}

\end{document}
